\scanpage{18}

$G$ un groupe plat et affine sur $S$, d'algèbre $B$. On suppose
\[
B=\varinjlim B_i,\qquad B_i\subset B\quad\text{sous-Modules avec}
\]

1) $B/B_i$ plat (donc $B_i\otimes B_i\subset B\otimes B$)

2) $\Delta B_i\subset B_i\otimes B_i$ donc $\check B_i$ est un faisceau d'algèbres associatives à unité

3) $B_i$ \deleted{projectif} \inserted{loc libre} de t.f. avec les $\check B_i$ idem.

$X$ un préschéma sur $S$, \deleted{muni de $\uline{L}$ \uline{ample} rel/$S$} \inserted{où $G$ opère.}

\deleted{\unsure{Ainsi}} On suppose que sur $X$ existe un faisceau \inserted{inv} $\uline{L}$ ample rel/$S$, sur lequel $G$ opère de façon compatible avec son action sur $X$.

Soit $E$ un $\underline O_X$-$G$-Module quasi-coh. \deleted{\unsure{on suppose que}}

1) Pour tt sous-module qu. coh $F$ de $E$, il existe un plus petit sous-module qu. coh $F'$ de $E$ \deleted{qui} stable par $G$ \unsure{contt} $F$. Sa \deleted{\unsure{définition}} formation commute à tt changement de base plat en $S$.
\inserted{Si $F$ est de t.f., de $=$ $F'$.}

\deleted{2) Si $S$ qu. cpt qu. sép., alors $E$ est $\varinjlim$ de sous-modules quasi-coh. de t.f. qui sont stables par $G$.}
\margin{\deleted{t.f.}}

\mspage{17}
